\chapter{Cuestionario II}

\section{Describir los procedimientos para configurar el ventilador pulmonar y hacer el registro de datos en una memoria USB}

La configuración del ventilador pulmonar Drager Babylog 8000+ comienza con la puesta en marcha del equipo y la verificación del autodiagnóstico inicial (niveles de oxigeno y aire), el cual permite comprobar el correcto funcionamiento neumático y físico del equipo, seguidamente para poder registrar los datos se realizan las configuraciones para establecer el flujo, presión y volumen, así como los tiempos de inspiración y expiración del paciente (frecuencia de respiración).

Una vez que el ventilador se encuentra correctamente configurado y operando, el sistema permite registrar y exportar información clínica y operativa del equipo. Para ello, el usuario accede a las funciones de monitorización o registro del ventilador, donde se almacenan datos como valores medidos, parámetros configurados y tendencias respiratorias del paciente. Luego se inserta una memoria USB en el puerto de datos del equipo y se selecciona la opción de exportación o transmisión de datos desde el menú correspondiente. El ventilador copia los registros almacenados en la memoria externa, generando archivos que posteriormente pueden ser analizados en un computador mediante programas de texto o hojas de cálculo, permitiendo conservar un historial sobre el funcionamiento del ventilador y de las variables respiratorias del paciente para fines clínicos, de investigación o documentación médica.

En forma de lista los pasos a seguir se detallan: 
\begin{enumerate}
	\item Encender el ventilador y verificar que el equipo complete su autodiagnóstico inicial.
	\item Acceder al menú principal de configuración del equipo desde el panel de control.
	\item Ajustar los parámetros de ventilación requeridos para el paciente (modo ventilatorio, presión inspiratoria, frecuencia respiratoria, FiO$_2$, relación I:E, entre otros).
	\item Verificar la correcta conexión del circuito respiratorio, sensores de flujo y sensores de oxígeno.
	\item Acceder al menú de monitorización y registro de datos del ventilador.
	\item Insertar una memoria USB en el puerto correspondiente del ventilador o en el módulo de interfaz de datos.
	\item Seleccionar la opción de exportación o registro de datos desde el menú de configuración o tendencias del sistema.
	\item Confirmar la operación de exportación para que el ventilador copie los datos registrados a la memoria USB.
	\item Esperar la confirmación en pantalla de que la transferencia de datos ha finalizado correctamente.
\end{enumerate}


